% !TeX root = ../main.tex
\subsection{Model Choice and Mathematical Rationale}
\label{sec:model-choice}

A suitable strength model must accommodate sparse, irregularly spaced results, update both opponents, distinguish draws, and quantify uncertainty before each match. It must also distinguish team ability from the advantage associated with playing at home. Elo-derived football predictors separate ratings from outcome regression~\cite{hvattum2010elo}, but a fixed-step point rating alone does not quantify uncertainty about ability. We use an internal fixed-step baseline for comparison. In our model, the latent strength difference is defined on a scale fixed by performance-noise variance and is related nonlinearly to outcome probability; it is not an observed physical measure of football quality.

Dynamic paired comparisons~\cite{glickman1999dynamic} distinguish latent ability from uncertainty in its estimation. Gaussian evolution provides a tractable starting point: inactivity increases uncertainty without changing expected strength. A smooth diffusion alone may, however, adapt slowly to persistent abrupt changes. We therefore consider both the smooth special case and a model with persistent jumps. The latter was selected using the development criteria described below. This selection is conditional on those criteria; a simpler model can have lower error in a stationary simulation setting.

TrueSkill~\cite{herbrich2006trueskill} incorporates uncertain skill, noisy performance, and a draw region, all relevant to an archive of home wins, draws, and away wins. Its application to game ratings does not establish its validity for national-team football. The time evolution, venue term, robust likelihood, and chronological checks used here require separate specification. Expectation propagation~\cite{minka2001ep} approximates inference across correlated latent states. Its use does not imply that the full posterior is Gaussian or that sequential inference is exact.

Goal-count models, such as that of Dixon and Coles~\cite{dixon1997scores}, could use additional score information. We do not adopt one here because the principal target is a home win and the recorded scores are not uniformly restricted to regulation time. Such an extension would require separate observation assumptions and evaluation, with goals and the resulting match outcome treated as dependent evidence. FIFA rankings and published Elo ratings were not evaluated as external benchmarks. Our results consequently provide no comparison with those systems.

Recent international-tournament forecasting also combines statistical learners. Groll et al.~\cite{groll2024combined} used generalized linear, random-forest, and extreme-gradient-boosting models to forecast goal counts for UEFA EURO 2024, incorporating ranking-derived predictors and additional team covariates. The present study instead focuses on historical context and uncertain pre-match strength using a result-based archive. Flexible ensembles provide a relevant alternative, but their predictive advantage over this implementation has not been tested.

The remaining methods address distinct inferential tasks. Direct standardization compares observed rates under a common match composition. Proper scoring rules~\cite{gneiting2007scoring} evaluate probabilities assigned to subsequent outcomes. Simulation-based uncertainty analysis evaluates comparison procedures under specified generating mechanisms. These methods complement strength estimation, but model-conditional interval coverage alone cannot establish predictive superiority in real matches.
